% Work-MAT | 2. rocnik | pripravovana hodina 25
% Návaznost: uživatelské připravy.zip (hodiny 01–20), autorská učebnice na GitHubu.
% Realisticky.cz: didaktická inspirace v posloupnosti a typech úloh; zadání jsou vlastní.
\documentclass[11pt,a4paper]{article}
\input{../../sablony/pripravy/priprava.sty}
\nastavPripravu{2. ročník}{Exponenciální rovnice}{25}{Vytýkání a substituce – mix}
\begin{document}
\blok{Motivace a návaznost}
Dnes nediktujeme metodu. Žáci mají z tvaru úlohy sami poznat, zda je rychlejší vytýkání společné mocniny, nebo substituce.
\poznamka{Orientační nabídka pro 45--55 minut: krátké zopakování a motivace, nový výklad, tři společné vzorové úlohy, samostatná práce. Vyber podle tempa třídy; zbývající zadání slouží jako rezerva.}
\blok{Výklad}
Pokud máme součet posunutých mocnin $a^{x+k}$, zkusíme vytknout $a^x$. Pokud máme $a^{2x}$ a $a^x$, často volíme kvadratickou substituci.
\blok{Společné příklady}
\begin{prikladReseniBox}[title={Příklad 1}]
Vyřeš, resp. urči: $3^{x+1}+3^x=108$.
\tcblower
\[\begin{aligned}3^x(3+1)=108 \\ 3^x=27=3^3 \\ x=3\end{aligned}\]
\end{prikladReseniBox}
\begin{prikladReseniBox}[title={Příklad 2}]
Vyřeš, resp. urči: $9^x-4\cdot3^x+3=0$.
\tcblower
\[\begin{aligned}t=3^x>0,\quad t^2-4t+3=0 \\ (t-1)(t-3)=0 \\ K=\{0,1\}\end{aligned}\]
\end{prikladReseniBox}
\begin{prikladReseniBox}[title={Příklad 3}]
Vyřeš, resp. urči: $2^x+2^{3-x}=6$.
\tcblower
\[\begin{aligned}t=2^x>0,\quad t+\frac{8}{t}=6 \\ K=\{1,2\}\end{aligned}\]
\end{prikladReseniBox}
\poznamka{Metody záměrně střídej. Po bloku vytýkání a substituce patří za DÚ právě tento smíšený soubor; opravuj volbu metody, nejen kořeny.}
\blok{Samostatné procvičování}
\begin{zadaniReseniBox}[title={Samostatná práce 1}]
\begin{enumerate}[label=\alph*)]
\item $2^{x+1}+2^x=24$.
\item $4^x-5\cdot2^x+4=0$.
\item $5^{x+1}-5^x=100$.
\item $9^x-10\cdot3^x+9=0$.
\item $3^{x+1}+3^x=108$.
\item $2^x+2^{3-x}=6$.
\end{enumerate}
\tcblower
\begin{enumerate}[label=\alph*)]
\item vytýkání, x=3.
\item substituce, K=\{0,2\}.
\item vytýkání, x=2.
\item substituce, K=\{0,2\}.
\item vytýkání, x=3.
\item substituce, K=\{1,2\}.
\end{enumerate}
\end{zadaniReseniBox}
\blok{Dobrovolné rozšíření}
\blok{Úkol na doma}
Při početních úlohách napiš postup, zvolenou metodu a zkoušku v původním zadání. U soustav kontroluj obě rovnice; u nerovnic výsledek zakresli na osu.
\begin{zadaniReseniBox}[title={Domácí úkol}]
\begin{enumerate}[label=\alph*)]
\item $3^{x+1}+3^x=108$.
\item $4^x-5\cdot2^x+4=0$.
\item $2^{x+2}-2^x=12$.
\item $9^x-4\cdot3^x+3=0$.
\item $3^{x+2}-3^x=72$.
\item $5^x+5^{3-x}=30$.
\item Vysvětli rozdíl mezi vytýkáním a substitucí..
\end{enumerate}
\tcblower
\begin{enumerate}[label=\alph*)]
\item vytýkání, x=3.
\item substituce, K=\{0,2\}.
\item vytýkání, x=2.
\item substituce, K=\{0,1\}.
\item vytýkání, x=2.
\item substituce, K=\{1,2\}.
\item Vytýkání izoluje společný činitel; substituce nahrazuje opakovaný výraz..
\end{enumerate}
\end{zadaniReseniBox}
\poznamka{Domácí úkol obsahuje 7 položek; předpokládaná domácí práce přibližně 15--30 minut. Po větších blocích jsou vědomě promíchané dříve probrané typy.}
\end{document}
